\section{Thành phần bổ sung}\label{sec:7}


Hai thành phần dưới đây chỉ làm nếu còn thời gian: \textbf{DMA theo chuỗi descriptor} (CPU ra lệnh một lần cho nhiều phép GEMM) và \textbf{scratchpad ping-pong} (nạp, tính, ghi của các phép GEMM kế nhau chạy chồng lên nhau). Chúng đã được đặc tả đầy đủ ở đây và dùng các offset, bit đã dành sẵn, nên khi hiện thực không phải sửa các mục 1 – 6.


\subsection{Nguyên tắc chung}\label{sec:7.1}


\paragraph{Nguyên tắc của phần bổ sung}\label{tab:7-1}
\begin{itemize}
\item Không thêm, bớt hay đổi port top-level (mục~\ref{sec:6}); không đổi offset hay bit của register map. Chỉ dùng CTRL[2], STATUS[3], các offset 0x40 … 0x48 và bit \sig{HW_CONFIG[17:16]}
\item Bật bằng tham số \sig{HAS_CHAIN}, \sig{HAS_PINGPONG}; phần cứng báo qua \sig{HW_CONFIG}. Firmware phải đọc \sig{HW_CONFIG} trước khi dùng phần bổ sung
\item Tác vụ START (chế độ đơn) giữ nguyên mọi hành vi của mục 1 – 6, dù phần bổ sung có hay không
\item Phần bổ sung chỉ thay đổi port nội bộ: DMA tách thành kênh đọc và kênh ghi chạy độc lập, thêm lệnh FETCH; các cổng SRAM thêm bit chọn bank; FSM thêm bộ lập lịch và hàng đợi descriptor
\end{itemize}


\subsection{DMA theo chuỗi descriptor}\label{sec:7.2}

CPU dựng sẵn một danh sách descriptor trong DRAM, mỗi descriptor mô tả một phép GEMM. CPU ghi địa chỉ descriptor đầu vào \sig{DESC_PTR}, thực hiện \sig{fence rw,rw}, ghi CTRL[2] = 1 (START\_CHAIN), rồi chờ DONE một lần cho cả chuỗi. Accelerator tự đọc từng descriptor và chạy hết chuỗi.



\begin{spectable}{|L{0.1}|L{0.08}|L{0.17}|L{0.12}|L{0.53}|}{Bố cục descriptor (16 word, 64 byte)}{tab:7-2}{\textbf{Offset} & \textbf{Word} & \textbf{Tên} & \textbf{Bit} & \textbf{Nội dung}}
0x00 & 0 & \sig{DIM} & [31:0] & M, N, K, bố cục như \sig{MATRIX_DIM} \\\hline

0x04 & 1 & \sig{A_ADDR} & [31:0] & Địa chỉ byte của A[0][0] \\\hline

0x08 & 2 & \sig{B_ADDR} & [31:0] & Địa chỉ byte của B[0][0]; bỏ qua khi \sig{REUSE_B} = 1 \\\hline

0x0C & 3 & \sig{C_ADDR} & [31:0] & Địa chỉ byte của C[0][0] hoặc Y[0][0] \\\hline

0x10 & 4 & \sig{LDA}, \sig{LDB} & [15:0], [31:16] & Stride theo hàng của A và B (byte); 0 = các hàng liền nhau \\\hline

0x14 & 5 & \sig{LDC} & [15:0] & Stride của C hoặc Y (byte); 0 = liền nhau; [31:16] dự trữ \\\hline

0x18 & 6 & \sig{FLAGS} & [4:0] & [0] \sig{LAST}, [1] \sig{WAIT_PREV}, [2] \sig{REUSE_B}, [3] \sig{POST_EN}, [4] \sig{RELU_EN}; [31:5] dự trữ \\\hline

0x1C & 7 & \sig{NEXT_PTR} & [31:0] & Địa chỉ descriptor kế tiếp, căn 64 byte; bỏ qua khi \sig{LAST} = 1 \\\hline

0x20 & 8 & \sig{BIAS_ADDR} & [31:0] & Như thanh ghi \sig{BIAS_ADDR}; chỉ dùng khi \sig{POST_EN} = 1 \\\hline

0x24 & 9 & \sig{QM0} & [31:0] & Như \sig{QUANT_M0}; chỉ dùng khi \sig{POST_EN} = 1 \\\hline

0x28 & 10 & \sig{QSHIFT} & [4:0] & Như \sig{QUANT_SHIFT}; [31:5] dự trữ \\\hline

0x2C & 11 & \sig{TAG} & [31:0] & Tùy phần mềm; phần cứng không đọc \\\hline

0x30–0x3C & 12–15 & — & — & Dự trữ, phải bằng 0 \\\hline

\end{spectable}

Với stride, địa chỉ phần tử là: A[i][k] tại \sig{A_ADDR} + i·\sig{LDA} + k; B[k][j] tại \sig{B_ADDR} + k·\sig{LDB} + j; C[i][j] tại \sig{C_ADDR} + i·\sig{LDC} + 4j; Y[i][j] tại \sig{C_ADDR} + i·\sig{LDC} + j. Stride = 0 được hiểu là đúng độ dài hàng.


\paragraph{Chế độ chuỗi}\label{tab:7-3}
\begin{itemize}
\item \sig{DESC_PTR} và mọi \sig{NEXT_PTR} được dùng phải căn 64 byte; nếu không, \sig{ERR_CODE} = 1. Mỗi descriptor được đọc đúng một lần, theo thứ tự chuỗi, bằng một burst INCR 16 beat (lệnh DMA FETCH)
\item Mỗi descriptor chạy các pha FETCH → kiểm tra → LOAD\_BIAS (nếu \sig{POST_EN}) → LOAD\_A → LOAD\_B (trừ khi \sig{REUSE_B}) → COMPUTE → STORE\_C, với nội dung giống hệt một tác vụ chế độ đơn có cùng tham số
\item Descriptor hợp lệ khi: M, N, K đúng miền; \sig{DIM[31:24]} = 0; mọi địa chỉ được dùng căn 4 byte; mỗi stride bằng 0 hoặc ≥ độ dài hàng (K cho A; N cho B; 4N cho C; N cho Y) và là bội của 4; bit và word dự trữ bằng 0; \sig{RELU_EN} kéo theo \sig{POST_EN}; \sig{REUSE_B} chỉ khi không phải descriptor đầu và K, N bằng của descriptor trước. Sai → \sig{ERR_CODE} = 1
\item \sig{REUSE_B} = 1: không nạp B; COMPUTE dùng B của descriptor ngay trước
\item \sig{WAIT_PREV} = 1: các pha LOAD của descriptor chỉ bắt đầu khi STORE\_C của mọi descriptor trước nó đã xong. Phần mềm phải đặt cờ này khi descriptor đọc dữ liệu do descriptor trước ghi ra
\item Chuỗi kết thúc ở descriptor đầu tiên có \sig{LAST} = 1. DONE = 1 khi STORE\_C của descriptor đó nhận phản hồi B cuối là OKAY
\item \sig{DESC_DONE_CNT} tăng 1 mỗi khi STORE\_C của một descriptor xong OKAY; các descriptor xong đúng theo thứ tự chuỗi
\item Lỗi: gọi f là descriptor có thứ tự nhỏ nhất bị lỗi. Không pha nào của descriptor thứ ≥ f được bắt đầu; các descriptor đứng trước f được chạy xong. Khi mọi khối đã rảnh: ERROR = 1, \sig{ERR_CODE} theo bảng mã lỗi (lỗi FETCH là mã 4), \sig{ERR_DESC_ADDR} = địa chỉ của f, \sig{DESC_DONE_CNT} = số thứ tự của f
\item Khi dùng stride lớn hơn độ dài hàng, DMA cắt burst theo từng hàng: burst không bao giờ trải qua khoảng hở giữa hai hàng, và các byte trong khoảng hở không bị đọc hay ghi
\item Accelerator không ghi vào descriptor; phần mềm không sửa descriptor và vùng dữ liệu của chuỗi đang chạy. Vùng đích của mọi descriptor không được chồng lên vùng descriptor
\item Ở chế độ chuỗi, tham số VPU lấy từ descriptor; các thanh ghi \sig{MATRIX_DIM}, \sig{BASE_ADDR_*}, \sig{POST_CTRL}, \sig{BIAS_ADDR}, \sig{QUANT_*} không được dùng
\end{itemize}


\subsection{Scratchpad ping-pong}\label{sec:7.3}

Mỗi buffer A, B, C, BIAS có hai bank. DMA kênh đọc nạp descriptor d+1 vào một bank trong lúc array tính descriptor d trên bank kia, và DMA kênh ghi chép kết quả của descriptor d−1 ra DRAM. Thời gian mỗi descriptor tiến tới max(nạp, tính, ghi) thay vì tổng của chúng. Phần này đòi hỏi chế độ chuỗi.



\begin{spectable}{|L{0.12}|L{0.88}|}{Chọn bank cho descriptor d}{tab:7-4}{\textbf{Buffer} & \textbf{Quy tắc}}
A, C & Đảo bank sau mỗi descriptor \\\hline

B & Đảo sau mỗi lần có LOAD\_B; với \sig{REUSE_B} = 1, dùng lại bank của descriptor trước \\\hline

BIAS & Đảo sau mỗi descriptor có \sig{POST_EN} = 1 \\\hline

\end{spectable}


\paragraph{Ping-pong}\label{tab:7-5}
\begin{itemize}
\item Ba khối thực thi: Đọc (FETCH, kiểm tra, LOAD\_*; kênh đọc DMA), Tính (COMPUTE; array và VPU), Ghi (STORE\_C; kênh ghi DMA). Mỗi khối làm một việc một lúc, theo thứ tự descriptor. Kênh đọc và kênh ghi được phép có burst AXI cùng lúc, mỗi kênh tối đa 1 burst outstanding
\item Khối Đọc làm lần lượt FETCH(d), kiểm tra d, các LOAD của d, rồi mới sang FETCH(d+1)
\item LOAD vào một bank chỉ bắt đầu khi COMPUTE cuối cùng đọc bank đó đã xong
\item COMPUTE(d) bắt đầu khi mọi LOAD của d và COMPUTE(d−1) đã xong, và STORE\_C cuối cùng đọc bank C của d đã xong
\item STORE\_C(d) bắt đầu khi COMPUTE(d) và STORE\_C(d−1) đã xong; do đó kết quả ra DRAM đúng theo thứ tự descriptor
\item \sig{WAIT_PREV}  chặn các LOAD, không chặn FETCH và kiểm tra
\item Tại mỗi thời điểm, mỗi bank có tối đa một bên truy cập. Hai bank khác nhau của cùng buffer được truy cập độc lập trong cùng chu kỳ; mỗi bank vẫn chỉ cần SRAM một cổng
\item Các điều kiện thứ tự và quyền truy cập bank nêu trên là điều kiện cần; phần cứng được phép bắt đầu một pha muộn hơn, không bao giờ sớm hơn. Nội dung DRAM sau tác vụ phải giống hệt khi chạy không ping-pong
\item Soft reset: abort cả hai kênh DMA đang bận; \sig{soft_clr} phát khi cả hai kênh đã rảnh
\item \sig{PERF_LOAD} + \sig{PERF_COMPUTE} + \sig{PERF_STORE} được phép lớn hơn \sig{PERF_TOTAL}; phần chênh là thời gian chạy chồng
\end{itemize}
